\chapter{METODOLOGIA}
\label{cap:metodologia}

\section{Classificação e Abordagem da Pesquisa}
\label{sec:classificacao-pesquisa}

Esta pesquisa classifica-se, quanto à sua natureza, como uma pesquisa \textbf{aplicada} de \textbf{desenvolvimento tecnológico} e inovação interdisciplinar, voltada à concepção e validação de produtos e soluções de impacto direto para o ambiente hospitalar e acadêmico da Universidade Federal de Sergipe (UFS). O trabalho articula as metodologias de Engenharia de Software às diretrizes científicas de segurança do paciente e gestão de riscos hospitalares.

O desenvolvimento do sistema estrutura-se em um ciclo iterativo de macroetapas rigorosas:
\begin{enumerate}
    \item \textbf{Identificação do Problema e Motivação}: Mapeamento das limitações dos sistemas passivos de notificação voluntária e das barreiras no ensino prático da metodologia IHI-GTT nos cursos de graduação em saúde;
    \item \textbf{Definição dos Objetivos da Solução}: Estabelecimento dos requisitos funcionais, ergonômicos e pedagógicos de uma plataforma computacional integrada;
    \item \textbf{Projeto e Desenvolvimento do Artefato}: Concepção da arquitetura desacoplada, modelos relacionais, estruturas flexíveis em JSONB e interface web de alta densidade informacional;
    \item \textbf{Demonstração e Simulação}: Aplicação do artefato em cenários simulados de casos clínicos e prontuários estruturados no ambiente acadêmico da UFS;
    \item \textbf{Avaliação e Homologação}: Mensuração de desempenho computacional, validação estrita do \textit{Quality Gate} (100\% de testes automatizados em linhas e branches) e avaliação de usabilidade heurística junto a discentes e docentes;
    \item \textbf{Comunicação e Disseminação}: Produção de documentação técnica, submissão de relatórios institucionais (parcial aos 6 meses e final aos 12 meses), elaboração de artigos científicos para periódicos indexados e depósitos de patentes junto ao INPI.
\end{enumerate}

A condução da pesquisa é intrinsecamente interdisciplinar, unindo os conhecimentos de Engenharia de Software e Sistemas Distribuídos do Departamento de Computação (DCOMP) à expertise clínica em segurança do paciente e auditoria hospitalar do Departamento de Enfermagem da UFS.

\section{Prospecção Tecnológica e Busca Sistemática de Anterioridade}
\label{sec:prospeccao-anterioridade}

Como elemento estruturante da metodologia de desenvolvimento e validação inventiva do SIGEA, realizou-se uma pesquisa sistemática de anterioridade patentária e prospecção tecnológica do estado da técnica, com levantamento concluído em \textbf{02 de setembro de 2026}. O objetivo primordial desta prospecção consistiu em mapear o panorama internacional de patentes e a literatura científica indexada, delimitando as lacunas existentes e comprovando a novidade, a atividade inventiva e a aplicação industrial da plataforma proposta.

O escopo do SIGEA reside no desenvolvimento de uma solução computacional para a identificação ativa e o rastreamento de eventos adversos (EAs) e incidentes clínicos a partir da metodologia \textit{Global Trigger Tool} (GTT) do \textit{Institute for Healthcare Improvement} (IHI). O principal diferencial inovador repousa em sua \textbf{aplicação pedagógica síncrona}: consolidar-se como um ambiente computacional simulado, interativo e não-punitivo para a formação clínica e tomada de decisão de graduandos em Enfermagem e Medicina. Ao estruturar a revisão retrospectiva de prontuários, a indexação determinística de gatilhos clínicos e a emissão de painéis gerenciais e relatórios sintético-analíticos em um ambiente controlado (\textit{sandbox}), a plataforma fomenta a transição de uma cultura hospitalar punitiva para uma abordagem sistemática de mitigação de riscos e melhoria contínua da qualidade assistencial.

\subsection{Metodologia e Metadados da Busca}
\label{subsec:metodologia-busca-anterioridade}

A formulação das estratégias de busca fundamentou-se nas diretrizes científicas de prospecção patentária e revisão sistemática da literatura. Foram consultadas bases referenciais nacionais e internacionais de patentes e literatura científica revisada por pares, garantindo cobertura multidisciplinar entre a Computação e as Ciências da Saúde.

A estruturação das expressões booleanas de busca articulou-se na interseção de três eixos temáticos fundamentais:
\begin{enumerate}
    \item \textit{Segurança do Paciente e Gestão de Riscos}: termos controlados e sinônimos como \textit{Patient Safety}, \textit{Adverse Events} e \textit{Risk Management};
    \item \textit{Engenharia de Software e Tecnologias de Informação}: termos como \textit{System}, \textit{Software}, \textit{Dashboard} e \textit{Application};
    \item \textit{Educação e Treinamento em Saúde}: termos direcionados à formação profissional, tais como \textit{Educational Tool}, \textit{Training}, \textit{Learning} e \textit{Teaching}.
\end{enumerate}

A Tabela \ref{tab:metadados-busca} consolida as bases consultadas, as tipologias documentais, as expressões booleanas parametrizadas e os quantitativos brutos e filtrados obtidos na prospecção.

\begin{table}[htbp]
\centering
\small
\caption{Metadados de busca e levantamento quantitativo por base de dados (02/09/2026)}
\label{tab:metadados-busca}
\begin{tabularx}{\textwidth}{>{\raggedright\arraybackslash}p{2.8cm} p{1.6cm} >{\raggedright\arraybackslash\scriptsize}X c}
\toprule
\textbf{Base de Dados} & \textbf{Tipologia} & \textbf{Estratégia e Expressão Booleana de Busca} & \textbf{Resultados} \\ 
\midrule
Scopus & Artigos & \texttt{TITLE-ABS-KEY(("adverse events" OR "patient safety" OR "risk management") AND ("educational tool" OR "training" OR "learning" OR "education" OR "teaching") AND ("system" OR "software" OR "dashboard" OR "application"))} & 26.185 \\ 
\midrule
Web of Science & Artigos & Tópico: \texttt{(("adverse events" OR "patient safety" OR "risk management") AND ("educational tool" OR "training" OR "learning" OR "education" OR "teaching") AND ("system" OR "software" OR "dashboard" OR "application"))} & 12.338 \\ 
\midrule
SciELO & Artigos & Tópico: \texttt{(("adverse events" OR "patient safety" OR "risk management") AND ("educational tool" OR "training" OR "learning" OR "education" OR "teaching") AND ("system" OR "software" OR "dashboard" OR "application"))} & 283 \\ 
\midrule
WIPO Patentscope & Patentes & \texttt{FP: ((("adverse events" OR "patient safety" OR "risk management") AND ("educational tool" OR "training" OR "learning" OR "education" OR "teaching") AND ("system" OR "software" OR "dashboard" OR "application")))} & 542 \newline (86 sob G16H) \\ 
\midrule
Portal CAPES & Artigos & Anyfield: \texttt{(("eventos adversos" OR "segurança do paciente" OR "gerenciamento de riscos") AND ("educação" OR "treinamento" OR "ensino") AND ("sistema" OR "software" OR "dashboard"))} & 292 \newline (65 filtrados) \\ 
\bottomrule
\end{tabularx}
\fonte{Elaborado pelo autor com base nos dados coletados em 02/09/2026.}
\end{table}

\subsection{Análise do Estado da Técnica em Patentes Globais}
\label{subsec:analise-patentes-globais}

Na base internacional de patentes \textbf{WIPO Patentscope}, a busca global recuperou inicialmente 542 documentos de patente. Ao refinar o conjunto para invenções indexadas sob o código IPC \textbf{G16H} (\textit{Tecnologia da Informação e Comunicação aplicada à Saúde}), obteve-se um universo relevante de 86 documentos com aderência temática direta. Dentre estes, destacam-se quatro documentos representativos do estado da técnica mais próximo, confrontados a seguir com as proposições do SIGEA:

\begin{enumerate}
    \item \textbf{Patente CN118888101A -- Chongqing Medical University} \cite{patente_cn118888101}:
    \begin{itemize}
        \item \textit{Título:} \textit{Patient safety training system and patient safety training method};
        \item \textit{Aspectos Técnicos:} Descreve um sistema computacional estruturado para capacitação de equipes e discentes de enfermagem em segurança clínica, mapeando vulnerabilidades assistenciais e correlacionando-as a matrizes teóricas de avaliação;
        \item \textit{Diferencial do SIGEA:} O documento chinês fundamenta-se na prescrição de módulos e currículos teóricos expositivos e estáticos. Em contraste, o SIGEA opera sobre uma abordagem empírica e dinâmica: o discente interage com simulações de auditoria retrospectiva baseada em prontuários fictícios, alimentando um \textit{dashboard} epidemiológico em tempo real, sem ancoragem em sequenciamento rígido de aulas teóricas.
    \end{itemize}

    \item \textbf{Patente CN111627533A -- Fortune Software Co., Ltd.} \cite{patente_cn111627533}:
    \begin{itemize}
        \item \textit{Título:} \textit{Hospital-wide adverse event active monitoring and management system};
        \item \textit{Aspectos Técnicos:} Propõe uma arquitetura de monitoramento ativo contínuo em rede hospitalar, dotada de módulos de extração de dados clínicos em produção e disparo automatizado de alertas de eventos sentinela;
        \item \textit{Diferencial do SIGEA:} Configura-se como uma plataforma puramente corporativa voltada ao faturamento e controle de risco administrativo intra-hospitalar. O SIGEA adota regras e fluxos análogos de triagem clínica em duas etapas (revisão primária e consenso), mas os reorienta a um ambiente controlado de aprendizagem e \textit{debriefing}, blindando o aluno de impactos assistenciais reais durante a curva de aprendizado.
    \end{itemize}

    \item \textbf{Patente CN113379384A -- Huaxun High Tech Co., Ltd.} \cite{patente_cn113379384}:
    \begin{itemize}
        \item \textit{Título:} \textit{Medical event big data platform};
        \item \textit{Aspectos Técnicos:} Trata-se de uma plataforma escalável de \textit{Big Data} para gestão do ciclo de vida de notificações assistenciais, análise de causa-raiz e visualização gerencial em redes de saúde;
        \item \textit{Diferencial do SIGEA:} O foco da anterioridade reside na infraestrutura de escalabilidade massiva de notificações hospitalares. O SIGEA transforma a notificação e o cálculo de taxas epidemiológicas (ex.: eventos adversos por 1.000 pacientes-dia) em objetos pedagógicos formativos, capacitando o discente antes de sua inserção nos cenários de estágio e residência clínica.
    \end{itemize}

    \item \textbf{Patente US20210074398A1 -- MedStar Health, Inc.} \cite{patente_us20210074398}:
    \begin{itemize}
        \item \textit{Título:} \textit{Evaluation of patient safety event reports from free-text descriptions};
        \item \textit{Aspectos Técnicos:} Emprega técnicas de Processamento de Linguagem Natural (PLN) e modelos de Aprendizado de Máquina para classificar relatos em texto livre de incidentes clínicos assistenciais;
        \item \textit{Diferencial do SIGEA:} A solução norte-americana propõe um classificador preditivo automatizado de risco corporativo. O SIGEA prioriza um buscador e rastreador determinístico de gatilhos operado ativamente pelo aluno, servindo como suporte à identificação de pistas clínicas antes de qualquer automação algorítmica profunda.
    \end{itemize}
\end{enumerate}

\subsection{Análise da Literatura Científica Correlata}
\label{subsec:analise-literatura-correlata}

Dentre as 65 publicações científicas revisadas por pares obtidas no Portal de Periódicos CAPES e SciELO com aderência ao escopo do projeto, três vertentes bibliográficas delimitam o estado da técnica científico e metodológico:

\begin{enumerate}
    \item \textbf{Validação Interdisciplinar de Softwares em Saúde}: Estudos recentes destacam a relevância de metodologias de validação de conteúdo e usabilidade de softwares médicos envolvendo comitês de juízes especialistas das áreas de Tecnologia da Informação e Enfermagem. Essa abordagem ratifica a necessidade de o SIGEA ser construído e validado de forma interdisciplinar entre o DCOMP e a Enfermagem da UFS, assegurando rigor computacional alinhado à ciência do cuidado seguro;
    \item \textbf{Monitoramento Eletrônico e Avaliação via Escala SUS}: Pesquisas direcionadas a sistemas de vigilância e rastreamento de incidentes hospitalares consolidam a aplicação da \textit{System Usability Scale} (SUS) para aferir métricas de usabilidade e carga cognitiva. O SIGEA adotará esse arcabouço metodológico na validação de aceitação com os discentes;
    \item \textbf{Incidentes Assistenciais e o Fenômeno da ``Segunda Vítima''}: A literatura científica aponta o elevado impacto psicológico e a sobrecarga emocional que acometem estudantes de graduação que vivenciam ou relatam eventos adversos em ambientes clínicos punitivos. O SIGEA atua preventivamente como um ambiente seguro e controlado onde a falha de triagem e a interpretação de gatilhos constituem oportunidades de reflexão guiada, mitigando a ansiedade antes da prática assistencial direta.
\end{enumerate}

\subsection{Enquadramento na Classificação Internacional de Patentes (IPC)}
\label{subsec:classificacao-ipc}

Considerando o escopo inventivo, as rotinas computacionais e a finalidade prática do software, o SIGEA classifica-se na interseção das seguintes classes da Classificação Internacional de Patentes (\textit{International Patent Classification} -- IPC):

\begin{enumerate}
    \item \textbf{G16H (Tecnologia da Informação e Comunicação aplicada à Saúde)}: Soluções digitais voltadas à gestão, diagnóstico, monitoramento e análise de dados clínicos hospitalares. Abrange o motor de rastreamento de gatilhos clínicos, a parametrização dos 6 módulos do IHI-GTT e os motores de cálculo de taxas epidemiológicas de danos hospitalares;
    \item \textbf{G09B (Sistemas e Instrumentos de Ensino e Treinamento)}: Dispositivos, matrizes lógicas e plataformas de software direcionados à educação profissional, instrução assistida por computador e simulação de cenários operacionais. Compreende o ambiente de turmas simuladas, resolução de prontuários fictícios em JSONB e correção formativa docente.
\end{enumerate}

\subsection{Avaliação de Patenteabilidade e Requisitos Legais de Inovação}
\label{subsec:requisitos-patenteabilidade}

Diante do estado da técnica nacional e internacional mapeado nesta busca de anterioridade, constata-se que o projeto do SIGEA atende plenamente aos requisitos legais de proteção tecnológica e propriedade intelectual estabelecidos pela Lei de Propriedade Industrial (Lei nº 9.279/1996) \cite{brasil1996lpi}:

\begin{itemize}
    \item \textbf{Novidade (Art. 11 da LPI)}: Não foram identificadas no estado da técnica anterioridades que combinem de forma intrínseca e síncrona um motor de rastreamento retrospectivo de gatilhos clínicos segundo a metodologia GTT com um ambiente computacional concebido para o treinamento pedagógico não-punitivo em nível de graduação em saúde;
    \item \textbf{Atividade Inventiva (Art. 13 da LPI)}: A arquitetura proposta não deriva de maneira óbvia ou trivial do estado da técnica para um especialista no assunto. A transposição da metodologia de auditoria retrospectiva do IHI em uma ferramenta interativa de apoio ao aprendizado e suporte à análise gerencial de causas representa uma fusão inovadora entre a Engenharia de Software e a Educação em Saúde;
    \item \textbf{Aplicação Industrial (Art. 15 da LPI)}: O sistema apresenta reprodutibilidade técnica, escalabilidade computacional em contêineres Docker e viabilidade imediata de implantação em hospitais universitários, centros de simulação clínica e matrizes curriculares acadêmicas.
\end{itemize}

A realização desta prospecção fundamenta formalmente o plano do projeto de proceder com o registro de programa de computador e depósito de privilégio de invenção junto ao Instituto Nacional da Propriedade Industrial (INPI), assegurando a proteção e a transferência tecnológica das inovações desenvolvidas na Universidade Federal de Sergipe.

\section{A Metodologia Global Trigger Tool (IHI-GTT)}
\label{sec:metodologia-gtt}

O \textit{Global Trigger Tool} baseia-se na identificação retrospectiva e ativa de eventos adversos através de uma lista finita de rastreadores ou ``gatilhos'' clínicos padronizados \cite{griffin2019ihi, pierdevara2020}. No SIGEA, a metodologia foi estruturada conforme as diretrizes canônicas internacionais:

\subsection{Módulos de Auditoria e Gatilhos Clínicos Padronizados}
O sistema contempla 6 módulos clínicos especializados, totalizando 53 gatilhos clínicos padronizados (Tabela \ref{tab:modulos-gtt}):
\begin{itemize}
    \item \textbf{Módulo Cuidados (C1 a C15)}: Rastreia eventos gerais da assistência no leito, como quedas com dano físico mensurável, lesões por pressão (LPP) em qualquer estágio, complicações de procedimentos invasivos, readmissões não planejadas em até 30 dias e infecções relacionadas à assistência à saúde (IRAS);
    \item \textbf{Módulo Medicação (M1 a M13)}: Rastreia incidentes farmacológicos críticos, incluindo hipoglicemia severa (glicemia $< 50$ mg/dL), alargamento crítico da coagulação (RNI $> 6$), elevação abrupta de escórias renais e administração de antídotos específicos como Naloxona, Flumazenil e Fitomenadiona (Vitamina K1);
    \item \textbf{Módulo Cirúrgico (S1 a S11)}: Focado no bloco cirúrgico e na Sala de Recuperação Pós-Anestésica (RPA), identificando retorno não programado à sala operatória, conversão de abordagem cirúrgica por intercorrência, lesão iatrogênica de órgãos e admissão não planejada em UTI;
    \item \textbf{Módulo Cuidados Intensivos / Terapia Intensiva (I1 a I4)}: Focado em pacientes críticos de UTI, rastreando pneumonia associada à ventilação mecânica (PAV), readmissão precoce em leito intensivo e extubações acidentais com reintubação de emergência;
    \item \textbf{Módulo Perinatal (P1 a P8)}: Rastreia complicações materno-fetais durante o parto e puerpério, como lacerações perineais graves de 3º e 4º graus e hemorragias puerperais anormais;
    \item \textbf{Módulo Serviço de Urgência / Pronto Atendimento (E1 e E2)}: Rastreia permanência excessiva em portas de entrada ($> 6$ horas) e retorno precoce ao pronto socorro em menos de 48 horas pós-alta.
\end{itemize}

\begin{table}[htbp]
\centering
\small
\caption{Distribuição dos Módulos Especializados e Quantidade de Gatilhos da Metodologia IHI-GTT}
\label{tab:modulos-gtt}
\begin{tabularx}{\textwidth}{l p{4.8cm} c X}
\toprule
\textbf{Código} & \textbf{Nome do Módulo} & \textbf{Gatilhos} & \textbf{Foco Assistencial} \\
\midrule
\texttt{CUIDADOS}  & Cuidados Assistenciais Gerais & 15 & Cuidados de enfermagem, integridade de pele e quedas \\
\texttt{MEDICACAO} & Medicamentos e Farmacoterapia & 13 & Fármacos de alta vigilância, toxicidade e antídotos \\
\texttt{CIRURGICO} & Cirúrgico e Pós-Operatório & 11 & Bloco cirúrgico, intraoperatório e pós-anestésico \\
\texttt{UTI}       & Cuidados Intensivos (UTI) & 4 & Suporte ventilatório, via aérea e sepse grave \\
\texttt{PERINATAL} & Perinatal e Obstétrico & 8 & Parto, puerpério e hemorragias obstétricas \\
\texttt{URGENCIA}  & Urgência e Emergência & 2 & Permanência prolongada e retorno não planejado \\
\midrule
\textbf{Total}     & \textbf{6 Módulos Especializados} & \textbf{53} & \textbf{Vigilância Clínica Hospitalar Integral} \\
\bottomrule
\end{tabularx}
\fonte{Adaptado de Griffin e Resar \cite{griffin2019ihi} e Pierdevara et al. \cite{pierdevara2020}.}
\end{table}

\subsection{Classificação de Dano segundo o NCC MERP}
Ao identificar a presença de um gatilho clínico no prontuário, o auditor avalia a existência de nexo causal e mensura a intensidade do dano através do índice padronizado do NCC MERP \cite{nccmerp2001}, categorizado em 9 níveis de gravidade (Tabela \ref{tab:ncc-merp}). O SIGEA implementa a regra formal de que exclusivamente as categorias de \textbf{E} a \textbf{I} constituem Eventos Adversos com dano real, enquanto de \textbf{A} a \textbf{D} configuram circunstâncias de risco ou incidentes sem dano.

\begin{table}[htbp]
\centering
\small
\caption{Categorização de Gravidade de Dano Conforme a Taxonomia Oficial NCC MERP}
\label{tab:ncc-merp}
\begin{tabularx}{\textwidth}{c c X}
\toprule
\textbf{Categoria} & \textbf{Dano Real?} & \textbf{Definição Operacional da Taxonomia} \\
\midrule
\textbf{A} & Não & Circunstâncias ou eventos com capacidade potencial para causar erro. \\
\textbf{B} & Não & O erro ocorreu, mas não atingiu o paciente (quase-falha ou \textit{near miss}). \\
\textbf{C} & Não & O erro atingiu o paciente, mas não provocou nenhum dano físico. \\
\textbf{D} & Não & O erro atingiu o paciente e exigiu monitoramento para confirmar a ausência de dano. \\
\textbf{E} & \textbf{Sim} & Dano temporário ao paciente com necessidade de intervenção terapêutica ativa. \\
\textbf{F} & \textbf{Sim} & Dano temporário com necessidade de iniciar ou prolongar a internação hospitalar. \\
\textbf{G} & \textbf{Sim} & Dano permanente ao paciente com sequela funcional irreversível. \\
\textbf{H} & \textbf{Sim} & Necessidade de intervenção de emergência para manutenção da vida (suporte vital). \\
\textbf{I} & \textbf{Sim} & Óbito do paciente atribuível direta ou indiretamente ao evento adverso. \\
\bottomrule
\end{tabularx}
\fonte{NCC MERP \cite{nccmerp2001}.}
\end{table}

\subsection{Métricas Epidemiológicas de Dano Hospitalar}
O sistema automatiza o cômputo dos três principais indicadores preconizados pelo IHI:
\begin{equation}
    \text{Taxa de EAs / 1.000 pt-dia} = \left( \frac{\sum \text{EAs Detectados (Cat. E--I)}}{\text{Total de Pacientes-Dia Auditados}} \right) \times 1.000
\end{equation}
\begin{equation}
    \text{\% de Admissões com EAs} = \left( \frac{\text{Prontuários com } \ge 1 \text{ EA}}{\text{Total de Prontuários Auditados}} \right) \times 100
\end{equation}
\begin{equation}
    \text{EAs / 100 Admissões} = \left( \frac{\text{Total de Eventos Adversos Detectados}}{\text{Total de Prontuários Auditados}} \right) \times 100
\end{equation}

\section{Ferramentas de Engenharia e Gestão da Qualidade Hospitalar}
\label{sec:ferramentas-qualidade}

Para que a plataforma transcenda a mera constatação de incidentes e atue como um instrumento didático e clínico de melhoria contínua, foram incorporadas ferramentas fundamentais da gestão da qualidade assistencial, alinhadas à prática da enfermagem \cite{sze2019, rodrigues2020}:

\begin{enumerate}
    \item \textbf{Diagrama de Ishikawa (Espinha de Peixe / 6M)}: Modela visualmente a análise de causa-raiz dos incidentes nas dimensões clássicas de Método, Máquina, Medida, Meio Ambiente, Mão de Obra e Material \cite{sze2019, rodrigues2020};
    \item \textbf{Ciclo PDCA / PDSA}: Estruturação iterativa nos quatro quadrantes de Planejar (\textit{Plan}), Fazer (\textit{Do}), Checar/Estudar (\textit{Check/Study}) e Agir (\textit{Act}) para formalização de melhorias e protocolos preventivos \cite{sze2019, rodrigues2020};
    \item \textbf{Matriz GUT}: Priorização quantitativa de riscos hospitalares baseada em notas de 1 a 5 para Gravidade ($G$), Urgência ($U$) e Tendência ($T$), calculando o índice de prioridade $\text{Escore} = G \times U \times T$ \cite{sze2019};
    \item \textbf{Matriz SWOT (FOFA)}: Avaliação diagnóstica situacional mapeando Forças (\textit{Strengths}), Oportunidades (\textit{Opportunities}), Fraquezas (\textit{Weaknesses}) e Ameaças (\textit{Threats}) na gestão do cuidado \cite{sze2019};
    \item \textbf{Método 5W2H / Matriz 5W3H}: Formalização executiva do plano de ação estruturado respondendo a O quê (\textit{What}), Por quê (\textit{Why}), Onde (\textit{Where}), Quando (\textit{When}), Quem (\textit{Who}), Como (\textit{How}) e Quanto Custa / Como Medir (\textit{How Much / How Many}) \cite{sze2019};
    \item \textbf{Brainstorming Estruturado}: Módulo interativo para sessões colaborativas de ideação livre e geração de hipóteses preventivas pela equipe assistencial \cite{sze2019, rodrigues2020}.
\end{enumerate}

\section{Arquitetura Computacional e Tecnologias Adotadas}
\label{sec:arquitetura-software}

O desenvolvimento do software é integralmente conduzido pelo discente Matheus Araujo Pereira no ambiente do \textbf{Antigravity IDE} executando nativamente no sistema operacional \textbf{Fedora 44 Workstation (x86\_64)}, sob a orientação dos professores Gilton (DCOMP) e Ana Waleska (Enfermagem).

Para garantir sustentabilidade e conformidade com as diretrizes do projeto, foram adotadas versões LTS (\textit{Long Term Support}) atuais de tecnologias abertas:

\subsection{Backend e Camada de Domínio}
\begin{itemize}
    \item \textbf{Linguagem}: Java 21 LTS (OpenJDK), usufruindo de recursos avançados como \textit{Pattern Matching}, tipos de registro (\textit{Records}) e alta eficiência de execução;
    \item \textbf{Framework Central}: Spring Boot 3.3.x com Spring Data JPA, Hibernate 6 e Bean Validation para integridade declarativa dos modelos;
    \item \textbf{Segurança e Autenticação}: Spring Security 6 com autenticação \textit{stateless} baseada em tokens JWT. O controle de autorização baseia-se no modelo RBAC (\textit{Role-Based Access Control}) contemplando exatamente três papéis: \texttt{ADMIN}, \texttt{PROFESSOR} e \texttt{STUDENT};
    \item \textbf{Mapeamento e Utilitários}: MapStruct para conversão rápida entre entidades JPA e DTOs, e Lombok para redução de código boilerplate;
    \item \textbf{Documentação da API}: SpringDoc OpenAPI 3 gerando especificação Swagger UI interativa em tempo real.
\end{itemize}

\subsection{Banco de Dados e Persistência Híbrida}
\begin{itemize}
    \item \textbf{SGBD}: PostgreSQL 16 LTS com suporte a chaves universais UUID (\texttt{uuid-ossp});
    \item \textbf{Versionamento de Esquema}: Flyway Migrations, garantindo reproducibilidade determinística da estrutura relacional;
    \item \textbf{Persistência Semi-Estruturada}: Utilização do tipo de dado nativo \texttt{JSONB} para armazenamento flexível de prontuários simulados complexos (evoluções, checagens de horários e exames) e históricos das ferramentas de qualidade.
\end{itemize}

\subsection{Frontend e Ergonomia Visual Clínica}
\begin{itemize}
    \item \textbf{Framework}: Angular 18+ baseado em \textit{Standalone Components}, diretivas nativas de controle de fluxo (\texttt{@if}, \texttt{@for}) e gerenciamento reativo de estado por \textit{Signals};
    \item \textbf{Design System Hospitalar}: TailwindCSS integrado a componentes do Angular Material, aplicando uma paleta de cores ergonômica: azuis clínicos profundos (\texttt{\#1E3A8A} e \texttt{\#2563EB}), verdes esmeralda suaves para confirmação (\texttt{\#059669}), cinzas de baixo cansaço visual (\texttt{\#F8FAFC}, \texttt{\#F1F5F9} e \texttt{\#0F172A}), âmbar (\texttt{\#D97706}) e carmim hospitalar para gravidades críticas (\texttt{\#DC2626});
    \item \textbf{Documentação Frontend}: Compodoc para mapeamento integral da arquitetura de componentes e serviços.
\end{itemize}

\section{Regras de Negócio e Governança de Usuários}
\label{sec:regras-negocio}

As seguintes regras foram estabelecidas como inegociáveis para a integridade do sistema:
\begin{itemize}
    \item \textbf{Domínio de E-mail Institucional Exclusivo}: Apenas endereços eletrônicos com o sufixo oficial \texttt{@academico.ufs.br} são aceitos. Qualquer outra tentativa de cadastro ou login é rejeitada no backend e frontend;
    \item \textbf{Níveis de Perfil}: O sistema restringe-se aos papéis \texttt{ADMIN}, \texttt{PROFESSOR} e \texttt{STUDENT};
    \item \textbf{Campo Matrícula}: Obrigatório \textbf{estritamente} para o perfil \texttt{STUDENT}. Para os perfis \texttt{ADMIN} e \texttt{PROFESSOR}, o campo é nulo e totalmente oculto na interface;
    \item \textbf{Credenciais e Primeiro Acesso}: O usuário é criado com uma senha provisória gerada automaticamente pelo sistema; no primeiro login, o sistema impõe o redirecionamento imediato para a tela de redefinição de senha. Ademais, o próprio usuário pode alterar sua senha a qualquer instante através de seu perfil pessoal;
    \item \textbf{Gestão de Acessos}: A criação, edição, ativação, inativação e exclusão de contas é prerrogativa exclusiva do perfil \texttt{ADMIN}.
\end{itemize}

\section{Modelagem Educacional e Gestão de Turmas}
\label{sec:modelagem-turmas}

O módulo educacional foi concebido para espelhar a estrutura didática dos cursos de graduação da UFS:
\begin{itemize}
    \item \textbf{Identificação da Turma}: Cadastro padronizado composto por: \textit{nome da matéria} + \textit{número da turma} + \textit{período letivo vigente} (ex.: \texttt{Física 3 - T06 - 2026.2});
    \item \textbf{Corpo Docente e Discente}: Cada turma possui exatamente um professor responsável (com possibilidade de substituição a qualquer momento pelo administrador, garantindo que a turma nunca fique desprovida de docente) e múltiplos discentes matriculados;
    \item \textbf{Atividades e Prazos}: O docente pode criar atividades avaliativas com prazos de entrega fixados. A atividade possui enunciado uniforme para a turma, contudo a submissão discente é individual;
    \item \textbf{Avaliação Pedagógica}: O docente corrige as submissões individuais atribuindo nota de 0.00 a 10.00 e fornecendo parecer qualitativo; não há limite máximo para a quantidade de atividades por turma;
    \item \textbf{Dashboards de Desempenho}: Painéis analíticos por turma para apoio pedagógico a professores e administradores;
    \item \textbf{Histórico Permanente}: Discentes e docentes mantêm acesso contínuo e irrestrito ao histórico de todas as turmas das quais já participaram;
    \item \textbf{Encerramento de Turma}: Operação exclusiva do administrador, condicionada à inexistência de atividades pendentes.
\end{itemize}

\section{Ergonomia, Resoluções Homologadas e Padrões de Usabilidade}
\label{sec:ergonomia-telas}

O sistema destina-se a estações de trabalho de secretarias acadêmicas e postos de enfermagem, focando na densidade informacional sem suporte mobile. O layout foi homologado nas seguintes resoluções de desktop:
\begin{itemize}
    \item \textbf{WUXGA}: $1920 \times 1200$ pixels;
    \item \textbf{HD+}: $1600 \times 900$ pixels;
    \item \textbf{HD (720p)}: $1280 \times 720$ pixels;
    \item \textbf{Ultrawide (21:9)}: $2560 \times 1080$ e $3440 \times 1440$ pixels.
\end{itemize}

Os padrões universais de usabilidade incluem: listagens padronizadas com exatamente 10 registros por página; paginação fluida; filtros de busca com \textit{debounce} de requisições; caixas modais para confirmação de ações irreversíveis; indicador visual global de carregamento (\textit{spinner}); e mensagens de retorno via notificações flutuantes (\textit{toasts}).

\section{Critério de Aceite Estrito (Quality Gate)}
\label{sec:quality-gate}

Para assegurar a confiabilidade do sistema, a transição entre fases do projeto obedece a um \textit{Quality Gate} inegociável:
\begin{enumerate}
    \item \textbf{Regra de Transição}: A Fase 2 só pode ser iniciada após a conclusão integral da Fase 1; similarmente, a Fase 3 depende da conclusão total da Fase 2;
    \item \textbf{Critérios Tripartites Obrigatórios}: Uma fase só é homologada quando atingir:
    \begin{itemize}
        \item 100\% de implementação funcional comprovada;
        \item 100\% de cobertura de testes automatizados (backend com JUnit 5, Mockito e Jacoco; frontend com Jest);
        \item 100\% de documentação de código gerada sem erros (OpenAPI/Javadoc e Compodoc).
    \end{itemize}
\end{enumerate}

\section{Aspectos Éticos, Aprovação Institucional e Simulação Total}
\label{sec:aspectos-eticos}

O projeto de pesquisa encontra-se formalmente aprovado pelo Comitê de Ética em Pesquisa da Universidade Federal de Sergipe (CEP/UFS), sob o Certificado de Apresentação de Apreciação Ética (CAAE) nº 91836925.8.0000.5546, e conta com parecer favorável emitido pela Gerência de Ensino e Pesquisa (GEP) do Hospital Universitário da UFS / EBSERH \cite{cep_ufs2025}.

Em observância à Lei Geral de Proteção de Dados Pessoais (LGPD, Lei nº 13.709/2018) e às resoluções do Conselho Nacional de Saúde (CNS nº 466/2012), a plataforma opera de maneira totalmente desacoplada e independente dos prontuários eletrônicos de produção do HU-UFS. Todos os casos clínicos, nomes, exames e históricos de eventos adversos monitorados (erros de medicação, broncoaspiração, flebites, extubações acidentais, lesões por pressão e infecções hospitalares) são \textbf{100\% simulados e fictícios} \cite{borzecki2020}, assegurando segurança ética irrestrita aos acadêmicos e preservando a privacidade dos pacientes reais. A aplicação pedagógica dessas simulações e a validação clínica da ferramenta são conduzidas pela equipe interdisciplinar e pelo corpo de usuários do Departamento de Enfermagem (DEN/UFS), discriminados no plano de trabalho (Seção \ref{sec:equipe-validacao}).
